% ================================================
% 文件: 06_传感器与换能器.tex
% 章节: 传感器与换能器
% 所属部分: 电子元器件与材料
% 版本: v1.0
% ================================================

\section{传感器与换能器}
\label{sec:sensors_transducers}

传感器是将非电信号转换为电信号的装置。
更严格地说，传感器是能感受规定的被测量并按一定规律将其转换为可用
输出信号的器件或装置，其输出通常是便于传输和处理的电量，如电压、
电流、电阻、电容或频率。与之相对，换能器泛指在两种能量形式之间
进行转换的器件，既包括由非电能转为电能的方向（如麦克风），也包括
由电能转为非电能的方向（如扬声器）。因此可以认为传感器是换能器中
以“获取信息”为目的的那一类，而扬声器一类以“输出能量”为目的的
器件则属于逆向换能器。在通信设备中，二者共同构成整机与外部物理
世界之间的接口。

\textbf{传感器的组成}：一个完整的传感器一般由敏感元件、转换元件和
信号调理电路三部分组成。敏感元件直接感受被测量并输出与之确定相关
的物理量；转换元件把该物理量变换为电参数；信号调理电路则把微弱的、
非标准的电参数放大、线性化并转换为标准输出。以电阻式压力传感器为
例，弹性膜片是敏感元件，贴附其上的应变片是转换元件，惠斯登电桥
与仪表放大器构成调理电路。有些传感器的敏感元件与转换元件合为一体，
如热电偶与压电晶体，它们直接把被测量转换为电动势或电荷，属于
自发电型（有源型）传感器，无需外部激励；而热敏电阻、应变片一类
则须外加激励电源才能输出信号，属于参量型（无源型）传感器。

\textbf{静态特性}：描述传感器在被测量缓变时输入输出关系的指标称为
静态特性，主要包括以下几项。灵敏度是输出增量与输入增量之比，
\[
S = \frac{\mathrm{d}y}{\mathrm{d}x}
\]
线性度（非线性误差）表示实际特性曲线偏离拟合直线的最大程度，
常以满量程输出的百分数表示，如 $\pm0.5\%$ FS。迟滞是正行程与
反行程特性曲线不重合的程度；重复性是同一条件下多次测量结果的
离散程度；分辨力是能引起输出可察觉变化的最小输入增量；量程是
允许测量的上下限之差；零点漂移与灵敏度漂移则描述输出随时间和
温度的缓慢变化。工程上通常把上述各项误差按平方和开方的方式
合成，得到传感器的综合精度等级。

\textbf{动态特性}：当被测量快速变化时，传感器的输出会滞后于输入。
多数传感器可近似为一阶系统，其阶跃响应为
\begin{equation}
y(t) = y_\infty \left(1 - e^{-t/\tau}\right)
\end{equation}
其中 $\tau$ 为时间常数，对应的上限截止频率为
$f_c = 1/(2\pi\tau)$。$\tau$ 越小则响应越快、可测频率越高。
温度传感器的时间常数主要由热容与换热系数决定，故测量流动气体
时应尽量减小敏感元件的体积并增强对流。压电式传感器则更接近
二阶系统，须使工作频率远低于其机械谐振频率，以避免共振引起的
幅值畸变。

\begin{itemize}
    \item \textbf{温度传感器}：如热敏电阻、热电偶
    \item \textbf{光传感器}：如光敏电阻、光电二极管
    \item \textbf{压力传感器}：如压电传感器
    \item \textbf{换能器}：如扬声器、麦克风
\end{itemize}

\textbf{温度传感器}：温度是最常被测量的物理量，通信设备中用它监视
功率放大器散热状况、补偿晶体振荡器频率漂移、控制风扇转速与
过热保护。热敏电阻是利用半导体材料电阻率随温度剧烈变化的原理
制成的，按温度系数符号分为负温度系数（NTC）与正温度系数（PTC）
两类。NTC 热敏电阻的阻值随温度升高而急剧下降，其阻温关系可用
指数式描述：
\begin{equation}
R_T = R_0\,\exp\!\left[B\left(\frac{1}{T} - \frac{1}{T_0}\right)\right]
\end{equation}
其中 $R_0$ 为参考温度 $T_0$（常取 $298\ \mathrm{K}$，即
$25\ ^{\circ}\mathrm{C}$）下的阻值，$B$ 为材料常数，通常在
$2000\sim5000\ \mathrm{K}$ 之间，温度均为绝对温度。NTC 的灵敏度
很高，但线性度差，须通过串并联电阻网络或查表法进行线性化；
它还常被串入电源回路作为浪涌抑制元件，上电瞬间因常温阻值大而
限制冲击电流，随后自身发热阻值下降而不影响正常工作。PTC 热敏
电阻在超过居里点后阻值突增，可作为自恢复保险与恒温加热元件。

金属热电阻是另一类温度传感器，以铂电阻（如 Pt100，即
$0\ ^{\circ}\mathrm{C}$ 时阻值为 $100\ \Omega$）最为典型，
其阻温关系在不太宽的范围内近似线性：
\[
R_t = R_0\left(1 + \alpha t\right)
\]
铂的电阻温度系数约为
$0.00385\ /^{\circ}\mathrm{C}$。铂电阻的精度、稳定性与互换性
远优于热敏电阻，但灵敏度低、成本高，且引线电阻会带来附加误差，
故精密测量须采用三线或四线制接法。

热电偶利用塞贝克效应工作：两种不同金属组成闭合回路，当两个
接点温度不同时回路中产生热电动势，其大小近似为
\begin{equation}
E = S\,(T_1 - T_2)
\end{equation}
$S$ 为塞贝克系数。常用的热电偶有 K 型（镍铬--镍硅）、
J 型（铁--康铜）、T 型（铜--康铜）等，K 型的塞贝克系数约为
$41\ \mu\mathrm{V}/^{\circ}\mathrm{C}$。由于输出取决于两端
温差，测量时必须已知参考端温度，工程上采用冰点槽或电子冷端
补偿电路来解决。热电偶的突出优点是测温范围极宽（可达
$1000\ ^{\circ}\mathrm{C}$ 以上）、结构简单、响应快，缺点是
输出电压极小，需要低漂移、高共模抑制比的前置放大器。此外，
半导体集成温度传感器利用 PN 结正向压降的温度系数（约
$-2\ \mathrm{mV}/^{\circ}\mathrm{C}$）实现测温，线性度好、
可直接输出数字量，在整机监控中应用最为方便。

\begin{table}[htbp]
\centering
\caption{常用温度传感器的性能对比}
\begin{tabular}{|p{2.2cm}|p{2.6cm}|p{2.4cm}|p{4.0cm}|}
\hline
\textbf{类型} & \textbf{测温范围} & \textbf{灵敏度/线性} & \textbf{特点与用途} \\
\hline
NTC 热敏电阻 & $-50\sim300\ ^{\circ}$C & 很高，非线性 & 廉价灵敏，需线性化；测温、浪涌抑制 \\
\hline
PTC 热敏电阻 & 居里点附近 & 阈值型 & 过流保护、自恢复保险、恒温加热 \\
\hline
铂电阻 Pt100 & $-200\sim650\ ^{\circ}$C & 中等，线性好 & 精度与稳定性高，需三/四线制接法 \\
\hline
热电偶（K 型） & $-200\sim1300\ ^{\circ}$C & 约 41 $\mu$V/$^{\circ}$C & 范围宽、响应快，需冷端补偿 \\
\hline
集成温度传感器 & $-55\sim150\ ^{\circ}$C & 线性，可数字输出 & 直接接入单片机，整机热管理 \\
\hline
\end{tabular}
\end{table}

\textbf{光传感器}：光传感器把光辐射转换为电信号，是光纤通信、
遥控接收、自动亮度调节与光电隔离的核心元件。光敏电阻（以硫化镉
CdS 为代表）利用光照产生的载流子降低半导体电阻率，暗电阻可达
兆欧级而亮电阻降至千欧级，灵敏度高但响应速度慢（毫秒至数十
毫秒量级），仅适用于照度检测一类的慢变场合。光电二极管工作在
反向偏置状态，入射光子在耗尽层产生电子空穴对形成光电流，其
大小与入射光功率成正比：
\begin{equation}
I_p = R_\lambda P_{opt}
\end{equation}
其中 $R_\lambda$ 为响应度，单位为 A/W。光电二极管线性好、
响应快（可达纳秒乃至皮秒量级），是光接收机的首选；PIN 结构
通过加宽本征层降低结电容以提高带宽，雪崩光电二极管（APD）
则利用倍增效应获得内部增益，用于长距离光纤通信。光电三极管
在光电二极管基础上附加放大作用，灵敏度提高 $\beta$ 倍，但线性
与带宽变差。光电池（太阳能电池）工作于零偏或正向状态，可直接
输出功率，属于自发电型器件。此外，把发光二极管与光电器件封装
在一起即构成光电耦合器，用于强弱电之间的电气隔离，是设备
抗干扰设计的常用手段。

\textbf{压力与力传感器}：压电式传感器利用某些晶体（石英、
钛酸钡、锆钛酸铅陶瓷）在受力时表面产生电荷的压电效应工作，
电荷量与作用力成正比：
\begin{equation}
Q = d\,F, \qquad U = \frac{Q}{C} = \frac{d\,F}{C}
\end{equation}
其中 $d$ 为压电常数，$C$ 为等效电容。由于电荷会通过绝缘电阻
缓慢泄漏，压电传感器不能测量静态力，只适于动态测量，且必须
配用输入阻抗极高的电荷放大器。压电效应还具有可逆性：外加电压
会使晶体产生形变，这就是逆压电效应，被用于制造石英晶体振荡器、
陶瓷滤波器、蜂鸣器与超声换能器。石英晶体的等效电路为串联
$L$--$C$--$R$ 支路与静态电容并联，其串联谐振频率
\[
f_s = \frac{1}{2\pi\sqrt{L_q C_q}}
\]
的稳定度可达 $10^{-6}$ 以上，因而成为通信设备频率基准的基础。

电阻应变式传感器则利用金属丝或半导体在受力伸长时电阻变化的
应变效应，其相对电阻变化与应变成正比：
\[
\frac{\Delta R}{R} = K\varepsilon
\]
$K$ 为灵敏系数，金属应变片约为 2，半导体应变片可达一百以上。
由于 $\Delta R/R$ 数量级仅为 $10^{-3}$ 甚至更小，必须接成
惠斯登电桥来提取信号。当四个桥臂全部为工作应变片且两两受
相反应变时（全桥接法），输出电压为
\begin{equation}
U_o = U_i\,\frac{\Delta R}{R} = U_i K \varepsilon
\end{equation}
若仅一臂为工作片（单臂接法），输出仅为上式的四分之一。
电桥接法不仅提高灵敏度，还能通过对称布片自动补偿温度引起的
共模变化。将应变片贴在弹性膜片或悬臂梁上，即构成压力、
称重与加速度传感器；采用集成工艺把扩散电阻直接做在硅膜片上，
则形成体积极小、灵敏度高的硅压阻式压力传感器。

\textbf{换能器：扬声器与麦克风}：扬声器是把电信号转换为声波的
换能器，以电动式（动圈式）最为常见。位于永久磁体气隙中的
音圈通过电流时受到安培力
\[
F = B I l
\]
其中 $B$ 为气隙磁感应强度，$l$ 为音圈在磁场中的有效长度。
该力驱动纸盆振动推动空气而发声。扬声器的主要参数包括额定
阻抗（常见 $4\ \Omega$、$8\ \Omega$、$16\ \Omega$，通信
耳机与话机常用更高阻抗）、额定功率、频率响应、灵敏度
（一米一瓦条件下的声压级）与谐振频率。需注意扬声器的阻抗
并非纯电阻，其模值随频率变化，在低频谐振点附近显著升高，
故功率放大器的负载匹配须按额定阻抗设计并留有余量。

麦克风（话筒）执行相反的转换。动圈式麦克风的结构与扬声器
互为逆过程，声波推动振膜带动线圈在磁场中运动，产生的感应
电动势为 $e = Blv$，其中 $v$ 为振膜速度；它无需供电、
坚固耐用、抗高声压，但灵敏度较低。电容式麦克风把振膜作为
电容的一个极板，声压引起间距变化从而改变电容量，在恒定
偏压下产生输出电压；其频率响应平直、灵敏度高，但需要极化
电压与前置放大。驻极体麦克风是电容式的改进，利用驻极体
材料自身的永久电荷取代外加极化电压，内置场效应管作阻抗
变换，只需一路直流偏置即可工作，体积小、成本低，是通信
终端与对讲设备中使用最广的话筒。传统碳粒式话筒则利用声压
改变碳粒堆的接触电阻，结构极简、输出电平高，曾长期用于
有线电话，但失真大、噪声高、频响窄。

\begin{table}[htbp]
\centering
\caption{常见传感器与换能器的转换关系}
\begin{tabular}{|p{2.6cm}|p{3.0cm}|p{2.6cm}|p{3.2cm}|}
\hline
\textbf{器件} & \textbf{转换关系} & \textbf{输出量} & \textbf{是否需外部激励} \\
\hline
NTC 热敏电阻 & 温度 $\rightarrow$ 电阻 & 电阻变化 & 需要（分压或电桥） \\
\hline
热电偶 & 温差 $\rightarrow$ 热电动势 & 毫伏级电压 & 不需要（自发电型） \\
\hline
光电二极管 & 光功率 $\rightarrow$ 光电流 & 电流 & 需反向偏压 \\
\hline
应变片 & 应变 $\rightarrow$ 电阻 & 电阻变化 & 需电桥激励 \\
\hline
压电元件 & 动态力 $\rightarrow$ 电荷 & 电荷/电压 & 不需要（自发电型） \\
\hline
动圈麦克风 & 声压 $\rightarrow$ 电动势 & 毫伏级电压 & 不需要 \\
\hline
驻极体麦克风 & 声压 $\rightarrow$ 电容变化 & 电压 & 需直流偏置 \\
\hline
扬声器 & 电流 $\rightarrow$ 声压 & 声波（逆向） & 由功放驱动 \\
\hline
\end{tabular}
\end{table}

\textbf{接口与信号调理}：传感器输出往往微弱且伴有干扰，必须经过
适当的调理电路才能送入后续处理。对于毫伏级差分输出（热电偶、
应变电桥），应采用高共模抑制比的仪表放大器，并注意屏蔽与
单点接地，信号电缆宜采用屏蔽双绞线且屏蔽层只在一端接地，以
避免地环流。对于高阻抗电荷输出（压电元件），须采用电荷放大器
或紧靠传感器安装的阻抗变换级，引线越短越好。对于电阻型
传感器，激励电流不宜过大，否则自身发热将引入误差。所有调理
电路都应设置与被测信号带宽相匹配的低通滤波，既抑制噪声，
又防止后续采样产生频谱混叠。

\textbf{选型与安装要点}：选择传感器时应依次考虑量程与过载能力、
精度与分辨力、响应速度、工作环境（温度、湿度、振动、
电磁干扰、腐蚀性介质）、输出形式与供电条件，以及长期稳定性
与标定周期。安装时须保证敏感元件与被测对象良好耦合：测温
元件应与被测面紧密接触并使用导热硅脂，避免受到辐射热与
气流的直接影响；测压元件应避免安装应力引入零点偏移；
声学器件应注意腔体密封与防止声反馈啸叫。此外，传感器属于
易受环境影响的元件，投入使用后应按规定周期进行校准，
用已知标准量核对其零点与灵敏度，必要时更新线性化系数。
正确理解各类传感器与换能器的原理、参数与局限，是保证通信
设备监控数据可信、音频通路清晰的基本条件。
